\begin{textsample}{POS Dim 1 – 2021\_c4 – Score -5.00 – 2021\_c4/enade\_2021\_pedag\_7adb919c15.txt}  \label{ex:f1_pos_379}
CAPÍTULO 4 PERCEPÇÃO DA PROVA Nas análises feitas neste capítulo, constam as percepções dos concluintes da Área de Pedagogia ( Licenciatura ) sobre a prova aplicada no Enade/2021.
Essas percepções foram mensuradas por meio de nove questões que avaliaram desde o grau de dificuldade da prova até o tempo gasto para concluí -la.
As percepções sobre a prova foram relacionadas ao desempenho dos estudantes, à Grande Região de funcionamento do curso, à Categoria Administrativa e ao tipo de Organização Acadêmica da IES.
O questionário de percepção da prova encontra -se ao final do Anexo VIII, onde está a reprodução do exame.
O desempenho dos estudantes foi classificado em quatro quartos.
Para tanto, esse desempenho foi ordenado de forma ascendente.
O percentil 25, P25, também conhecido como primeiro quartil, é a nota de desempenho que deixa um quarto ( 25 \% ) dos valores observados abaixo e três quartos acima.
Na Figura 1, apresenta -se uma ilustração desse conceito.
O quarto inferior de desempenho é composto pelas notas abaixo do primeiro quartil.
Já o percentil 75, P75, também conhecido como terceiro quartil, é o valor para o qual há três quartos ( 75 \% ) dos dados abaixo e um quarto acima dele.
O quarto superior de desempenho é composto pelas notas iguais ou acima do terceiro quartil.
O percentil 50, P50, também conhecido como mediana, é o valor que divide as notas em dois conjuntos de igual tamanho.
O segundo quarto inclui os valores entre o primeiro quartil ( P25 ) e a mediana.
O terceiro quarto contém os valores entre a mediana ( P50 ) e o terceiro quartil ( P75 ).
Vale ressaltar que percentis, quartis e medianas são pontos que, não obrigatoriamente, pertencem ao conjunto original de dados, ao passo que os quartos são subconjuntos dos dados originais.
A seguir, serão apresentados gráficos com resultados selecionados, relativos às nove questões avaliadas por grupos de estudantes.
Nas barras dos gráficos, apresenta -se o percentual de alunos que assinalaram uma das opções ou a soma das porcentagens daqueles que assinalaram duas ( ou três ) delas.
Por exemplo, para as questões 1 e 2, nos gráficos, é apresentada a porcentagem total de participantes que assinalaram as opções ( D ) Difícil e ( E ) Muito difícil.
Em cada barra, foram assinalados, também, os extremos do intervalo de confiança de 95 \% como linhas verticais unidas por uma linha horizontal na forma da letra H maiúscula.
O estimador de um parâmetro com um certo nível de confiança ( por exemplo, 95,0 \% ) deve conter o parâmetro no intervalo de confiança em 95 \% das vezes.
Na comparação entre os estimadores dos parâmetros de duas classes de uma dada categoria ( por exemplo : Norte e Nordeste nas Grandes Regiões, de primeiro e último quarto dentro de desempenho, de IES Públicas e Privadas, ou de Universidades e Faculdades ), associados aos seus respectivos intervalos de confiança, diz -se que não há diferença estatisticamente significativa entre os parâmetros das duas categorias quando há interseção entre os intervalos de confiança, e que há diferença, se os intervalos de confiança são disjuntos ( para maiores informações vide Glossário ).
Nas Tabelas do Anexo II, são apresentados os valores absolutos e a distribuição percentual18 das alternativas válidas das nove questões, segundo o mesmo recorte de desempenho dos alunos e da Grande Região de funcionamento do curso, de Categoria Administrativa e do tipo de Organização Acadêmica da IES. 18 Cumpre lembrar uma das convenções para tabelas numéricas ( ao final do Sumário ) sobre a possibilidade de a soma das partes não resultar em 100 \% por questões de arredondamento. 4.1 GRAU DE DIFICULDADE DA PROVA 4.1.1 COMPONENTE DE FORMAÇÃO GERAL Ao avaliarem “ Qual o grau de dificuldade desta prova na parte de Formação Geral? ” ( Questão 1 ), 40,7 \% do grupo de inscritos e presentes optaram pelas alternativas Difícil ou Muito difícil.
Entretanto, para mais da metade dos estudantes ( 53,6 \% ), o Componente de Formação Geral da prova foi avaliado como tendo grau de dificuldade Médio ( Gráfico 4.1, Gráfico 4.2 e, no Anexo II, a Tabela II.1 ).
O percentual de estudantes que consideraram a prova Difícil ou Muito difícil foi maior na região Sul, onde a proporção foi de 42,8 \%, enquanto a de menor proporção foi a Nordeste, com 35,8 \%.
No Gráfico 4.1, é possível observar que as diferenças entre as regiões Sul e Centro-Oeste com as regiões Norte, Nordeste e Sudeste são estatisticamente significativas, bem como entre as regiões Nordeste e Sudeste.
Nas Grandes Regiões, a proporção de presentes à prova que avaliaram o Componente de Formação Geral como tendo grau de dificuldade Médio esteve entre 52,2 \%, na região Sul, e 57,7 \%, na região Nordeste. presentes à prova que avaliaram o Componente de Formação Geral como tendo grau de dificuldade Médio esteve entre 52,2 \%, na região Sul, e 57,7 \%, na região Nordeste.
O percentual de alunos que consideraram o componente de Formação Geral da prova Difícil ou Muito difícil foi decrescente em função dos quartos de desempenho : 46,0 \%, no primeiro quarto, e 33,3 \%, no quarto, o grupo de melhor desempenho na prova.
Nos quartos de desempenho intermediários, a proporção de alunos que consideraram a prova Difícil ou Muito difícil foi de 43,3 \%, no segundo quarto, e de 40,5 \%, no terceiro quarto.
As diferenças entre todos os quartos de desempenho são estatisticamente significativas.
Para todos os quartos de desempenho, a alternativa modal para essa pergunta foi a Médio, com 47,9 \% e 59,6 \% dos respondentes nos quartos extremos, primeiro e quarto, respectivamente.
O percentual de estudantes que consideraram o componente de Formação Geral da prova Difícil ou Muito difícil foi maior entre os estudantes de CEFET/IF, alcançando 44,9 \%, enquanto entre os estudantes das Universidades, a Organização Acadêmica com a menor incidência, o percentual alcançou 40,5 \%.
No Gráfico 4.3, é possível observar que as diferenças entre CEFET/IF e as demais Organizações Acadêmicas são estatisticamente significativas.
Nas Organizações Acadêmicas, a proporção de presentes à prova que avaliaram o Componente de Formação Geral como de grau de dificuldade Médio esteve entre 50,4 \%, nos CEFET/IF, e 53,9 \%, nas Universidades ( ver também a Tabela II.2, no Anexo II ).
O percentual de estudantes que consideraram o componente de Formação Geral da prova Difícil ou Muito difícil foi maior entre os estudantes de cursos de IES Privadas ( 41,8 \% ), com diferença estatisticamente significativa da proporção registrada em IES Públicas ( 35,0 \% ).
Já a proporção de presentes à prova que avaliaram o Componente de Formação Geral como de grau de dificuldade Médio foi de 57,6 \%, entre os participantes de IES Públicas, e de 52,9 \%, entre os de IES Privadas ( ver Gráfico 4.4 ). 4.1.2 COMPONENTE DE CONHECIMENTO ESPECÍFICO Ao responderem à Questão 2 – “ Qual o grau de dificuldade desta prova na parte de Componente Específico? ” – 39,8 \% do grupo de estudantes classificaram -na como Difícil ou Muito difícil.
Além disso, o Componente modal de Conhecimento Específico da prova foi considerado como de grau de dificuldade Médio por 54,7 \% dos alunos ( Gráfico 4.5, Gráfico 4.6, e, no Anexo II, a Tabela II.3 ).
A análise das respostas dos estudantes quanto ao grau de dificuldade do Componente de Conhecimento Específico da prova, agregado por Grande Região, mostra diferenças estatisticamente significativas entre proporções de alunos que a avaliaram como Difícil ou Muito difícil.
Tais diferenças são observadas entre a região com a maior proporção, a Sul ( 41,5 \% ), e as regiões Norte ( 34,4 \% ), Nordeste ( 34,4 \% ), Sudeste ( 39,8 \% ) e Centro-Oeste ( 39,8 \% ), bem como entre as regiões Norte e Nordeste com as regiões Sudeste e Centro-Oeste.
O percentual de alunos que classificaram o grau de dificuldade como Médio, no Componente de Conhecimento Específico, variou de 53,4 \% a 58,6 \%, nas regiões Sul e Nordeste, respectivamente.
Considerando -se a avaliação da dificuldade das questões do Componente de Conhecimento Específico da prova de acordo com o desempenho dos estudantes, observa -se uma tendência decrescente dos que consideraram a parte específica Difícil ou Muito difícil, variando de 35,9 \% ( último quarto ) a 43,2 \% ( primeiro quarto ).
As demais proporções foram de 41,0 \% e 39,2 \%, respectivamente, no segundo e terceiro quartos.
Observa -se que há diferença estatisticamente significativa dos resultados entre todos os quartos de desempenho.
Já os que responderam que o grau de dificuldade das questões do Componente de Conhecimento Específico da prova foi Médio atingiram percentagens que variaram de 50,5 \%, no primeiro quarto, a 58,3 \%, no último.
Considerando -se a avaliação da dificuldade das questões do Componente de Conhecimento Específico da prova, o percentual de estudantes que as consideraram Difícil ou Muito difícil foi maior entre os estudantes de CEFET/IF, atingindo 43,3 \% ; já entre os de Faculdades, a Organização Acadêmica com a menor incidência, essa proporção foi de 39,0 \%.
No Gráfico 4.7, é possível observar que a diferença de Centros Universitários ( 40,7 \% ) e de CEFET/IF para Faculdades e Universidades ( 39,5 \% ) é estatisticamente significativa.
Nas Organizações Acadêmicas, a proporção de presentes à prova que consideraram o Componente de Conhecimento Específico como de grau de dificuldade Médio esteve entre 52,2 \%, nos CEFET/IF, e 55,0 \%, nas Faculdades ( ver também Tabela II.4, no Anexo II ).
O percentual de estudantes que consideraram o Componente Específico da prova Difícil ou Muito difícil foi maior entre os estudantes de cursos de IES Privadas ( 40,5 \% ), com diferença estatisticamente significativa da proporção registrada em IES Públicas ( 35,9 \% ).
Já a proporção de presentes à prova que avaliaram esse componente como de grau de dificuldade Médio foi de 57,3 \%, entre os participantes de IES Públicas, e de 54,2 \%, entre os de IES Privadas ( ver Gráfico 4.8 e Tabela II.4, no Anexo II ). 4.2 EXTENSÃO DA PROVA EM RELAÇÃO AO TEMPO TOTAL Indagados quanto à extensão da prova, em relação ao tempo total oferecido para a sua resolução ( Questão 3 ), os estudantes apontaram, com maior incidência, a alternativa que considerava a extensão Adequada, em todas as agregações consideradas ( Gráfico 4.9, Gráfico 4.10, e, no Anexo II, a Tabela II.5 ).
O percentual de alunos que responderam ser a extensão da prova Adequada foi de 46,7 \%.
Já 45,2 \% dos inscritos presentes consideraram que a prova foi Longa ou Muito longa, e 8,1 \% a avaliaram como Curta ou Muito curta.
Entre as Grandes Regiões, os que avaliaram a prova como Longa ou Muito longa, em relação ao tempo total destinado à sua resolução, atingiram uma porcentagem que variou de 36,6 \%, na região Norte, até 46,5 \%, na região Sudeste.
Há diferença estatisticamente significativa entre todas as regiões, com exceção apenas da diferença entre as regiões Nordeste e Sul.
Considerando -se o desempenho dos alunos, nota -se ainda que 49,1 \% consideraram a extensão da prova Adequada, no quarto de desempenho inferior, e 43,6 \%, no de melhor desempenho ( quarto superior ).
Nos quartos intermediários, essa proporção foi de 48,1 \%, no segundo quarto, e de 46,3 \%, no terceiro.
No Gráfico 4.10, observa -se que as proporções de estudantes que consideraram a prova Longa ou Muito longa apresentaram tendência crescente em função dos quartos de desempenho, sendo a menor proporção dessa resposta no primeiro quarto de desempenho ( 42,2 \% ), e a maior, no quarto superior ( 49,1 \% ).
Pode -se constatar que há diferença estatisticamente significativa entre todos os quartos de desempenho.
Considerando -se a extensão da prova, o percentual de estudantes que avaliaram a prova como Longa ou Muito longa, em relação ao tempo total destinado à sua resolução, foi maior entre os estudantes de CEFET/IF, alcançando 51,0 \%, enquanto entre os estudantes de Faculdades, a Organização Acadêmica com a menor incidência, esse índice atingiu 44,4 \%.
Entre os participantes de Centros Universitários, o índice atingiu 44,6 \%, e entre os de Universidades, 45,4 \%.
No Gráfico 4.11, é possível observar que só não há diferença estatisticamente significativa entre Centros Universitários e Faculdades.
A proporção de presentes à prova que consideraram a sua extensão Adequada, em relação ao tempo total destinado à sua resolução, esteve entre 41,2 \%, nos CEFET/IF, e 47,7 \%, nos Centros Universitários. ( ver também Tabela II.6, no Anexo II ).
O percentual de estudantes que consideraram a prova Longa ou Muito longa, em relação ao tempo total destinado à sua resolução, foi maior entre os estudantes de cursos de IES Públicas ( 49,7 \% ), havendo diferença estatisticamente significativa da proporção em relação aos de IES Privadas ( 44,3 \% ).
Já a proporção de presentes à prova que consideraram a extensão Adequada foi de 41,6 \%, entre os estudantes de IES Públicas, e de 47,7 \%, entre os de IES Privadas ( ver Gráfico 4.12 e Tabela II.6, no Anexo II ). 4.3 COMPREENSÃO DOS ENUNCIADOS DAS QUESTÕES 4.3.1 COMPONENTE DE FORMAÇÃO GERAL Em relação aos enunciados das questões do Componente de Formação Geral ( Questão 4 ), as opiniões foram positivas, já que 73,5 \% dos alunos avaliados consideraram todos ou a maioria dos enunciados das questões “ claros e objetivos ” ( Gráfico 4.13, Gráfico 4.14 e, no Anexo II, a Tabela II.7 ).
Na análise regional, a percentagem de estudantes que avaliaram que todos ou a maioria dos enunciados das questões do Componente de Formação Geral estavam “ claros e objetivos ” variou de 69,0 \%, na região Nordeste, a 75,0 \%, na região Centro-Oeste.
A análise das percepções dos estudantes sobre a clareza e a objetividade dos enunciados permite afirmar que todos ou a maioria dos enunciados de questões relativas ao Componente de Formação Geral, foram considerados “ claros e objetivos ” entre a maior parte dos respondentes ( maior ou igual a 69,0 \%, em todas as regiões, e maior ou igual a 65,3 \%, em todos os quartos de desempenho ).
Segundo o desempenho, observa -se que a proporção dos que emitiram essa opinião cresce, conforme o desempenho aumenta, com diferenças estatisticamente significativas entre todos os quartos de desempenho.
No quarto superior, a clareza e a objetividade de todos ou da maioria dos enunciados das questões foram percebidas por 81,4 \%.
A proporção dos presentes que consideraram que todos ou a maioria dos enunciados das questões do Componente de Formação Geral estavam “ claros e objetivos ” foi maior entre os estudantes de Universidades, chegando a 74,1 \%, enquanto entre os de Faculdades, a Organização Acadêmica com a menor incidência, esse índice atingiu 71,0 \%.
Entre os estudantes de Centros Universitários, o índice atingiu 73,4 \%, e entre os de CEFET/IF, 71,3 \%.
No Gráfico 4.15, é possível observar que as diferenças de Faculdades para Universidades e Centros Universitários são estatisticamente significativas ( ver também Tabela II.8, no Anexo II ).
O percentual de estudantes que consideraram que todos ou a maioria dos enunciados das questões do Componente de Formação Geral estavam “ claros e objetivos ” foi maior entre os estudantes de cursos de IES Públicas ( 74,2 \% ), e com diferença estatisticamente significativa da proporção em IES Privadas ( 73,4 \% ).
Um maior detalhamento pode ser encontrado no Gráfico 4.16 e na Tabela II.8, no Anexo II. 4.3.2 COMPONENTE DE CONHECIMENTO ESPECÍFICO Em relação aos enunciados das questões do Componente de Conhecimento Específico da prova, para 73,9 \% dos estudantes avaliados da Área de Pedagogia ( Licenciatura ), a clareza e a objetividade ( Questão 5 ) estavam presentes em todas ou na maioria das questões ( Gráfico 4.17 a Gráfico 4.20, e no Anexo II, as Tabelas II.9 e II.10 ).
A maioria dos estudantes de todas as Grandes Regiões brasileiras considerou “ claros e objetivos ” todos ou a maioria dos enunciados das questões do Componente de Conhecimento Específico da prova, percentual sempre maior ou igual a 69,7 \%.
A diferença entre a região Nordeste e todas as demais regiões é estatisticamente significativa, bem como entre a região Sudeste e as regiões Sul e Centro-Oeste e entre a região Norte e a região Sul.
A proporção de estudantes que consideraram os enunciados das questões “ claros e objetivos ” apresenta uma tendência crescente em relação ao aumento de desempenho : mais elevada no quarto superior ( 81,9 \% ), se comparada ao quarto inferior de desempenho ( 65,2 \% ).
As diferenças entre todos os quartos de desempenho são estatisticamente significativas.
A proporção dos presentes que consideraram que todos ou a maioria dos enunciados das questões do Componente de Conhecimento Específico estavam “ claros e objetivos ” foi maior entre os estudantes de Universidades, atingindo 74,4 \%.
Entre os estudantes de Faculdades e de CEFET/IF, ambas Organizações Acadêmicas com a menor incidência, essa proporção atingiu 71,4 \% e 71,3 \% respectivamente, sendo intermediária a proporção alcançada entre os estudantes de Centros Universitários ( 74,0 \% ).
No Gráfico 4.19, é possível observar que as diferenças de Universidades para Faculdades e CEFET/IF são estatisticamente significativas, bem como entre Centros Universitários e Faculdades ( ver também Tabela II.10, no Anexo II ).
O percentual de estudantes que consideraram que todos ou a maioria dos enunciados das questões do Componente de Conhecimento Específico estavam “ claros e objetivos ” foi maior entre os estudantes de cursos de IES Públicas ( 75,2 \% ), com diferença estatisticamente significativa da proporção registrada entre os de IES Privadas ( 73,7 \% ).
Ver Gráfico 4.20 e Tabela II.10, no Anexo II, para maior detalhamento. 4.4 SUFICIÊNCIA DAS INFORMAÇÕES/INSTRUÇÕES FORNECIDAS Ao se avaliarem as informações/instruções fornecidas para a resolução das questões ( Questão 6 ), constatou -se que 80,2 \% dos respondentes da Área de Pedagogia ( Licenciatura ) de todo o Brasil afirmaram que essas eram até excessivas ou suficientes em todas ou na maioria das questões ( Gráfico 4.21, Gráfico 4.22 e, no Anexo II, a Tabela II.11 ).
Quanto à distribuição de respondentes pelas Grandes Regiões, observa -se que a proporção de estudantes que consideraram as informações/instruções fornecidas até excessivas ou suficientes em todas ou na maioria das questões foi sempre superior ou igual a 77,4 \%, chegando a 80,8 \%, na região Norte.
Existem diferenças estatisticamente significativas entre a região Nordeste e todas as demais regiões.
Levando -se em conta o desempenho dos participantes, observa -se uma tendência crescente das proporções de participantes que avaliaram as informações/instruções como até excessivas ou suficientes em todas ou na maioria das questões, havendo diferenças estatisticamente significativas entre todos os quartos de desempenho, conforme mostrado no Gráfico 4.22.
O percentual foi mais elevado no quarto superior ( 87,9 \% ), percentual superior ao da média nacional ( 80,2 \% ).
Já no quarto inferior, a suficiência das informações/instruções declarada como até excessiva, em todas ou na maioria das questões foi percebida por 71,7 \% dos respondentes.
Considerando -se as informações/instruções fornecidas para a resolução das questões ( Questão 6 ), a proporção que considerou que essas eram até excessivas ou suficientes em todas ou na maioria das questões foi maior entre os estudantes de Centros Universitários, totalizando 81,0 \%.
As Faculdades foram a Organização Acadêmica em que houve a menor proporção : 78,3 \%.
No Gráfico 4.23, é possível observar que as diferenças entre Faculdades, Centros Universitários e Universidades ( 80,2 \% ), são estatisticamente significativas ( ver também Tabela II.12, no Anexo II ).
O percentual de estudantes que consideraram que as informações/instruções fornecidas para a resolução das questões eram até excessivas ou suficientes em todas ou na maioria das questões foi maior entre os estudantes de cursos de IES Públicas ( 82,5 \% ), havendo diferença estatisticamente significativa da proporção registrada entre os de IES Privadas ( 79,7 \% ) ( ver Gráfico 4.24 e Tabela II.12, no Anexo II ). 4.5 DIFICULDADE ENCONTRADA AO RESPONDER À PROVA Indagados sobre as dificuldades com as quais se depararam ao responder à prova ( Questão 7 ), 11,6 \% dos estudantes apontaram o Desconhecimento do conteúdo.
Para 49,9 \%, a Forma diferente de abordagem do conteúdo foi indicada como dificuldade.
Já a Falta de motivação para fazer a prova foi a dificuldade apontada por 13,7 \% dos respondentes.
Considerando -se todo o Brasil, 20,0 \% dos respondentes afirmaram que não tiveram qualquer tipo de dificuldade para responder à prova ( Tabelas II.13 e II.14, no Anexo II ).
Nos Gráficos 4.25 a 4.28, são apresentados os percentuais de estudantes que apontaram o Desconhecimento do conteúdo como “ dificuldade ao responder à prova ”.
Na análise por Grandes Regiões, o percentual de inscritos e presentes que apontaram o Desconhecimento do conteúdo como “ dificuldade ao responder à prova ” não superou 12,5 \%.
Os percentuais variaram de 10,1 \%, na região Nordeste, a 12,5 \%, na região Sul.
A Forma diferente de abordagem do conteúdo foi a escolha modal dos estudantes, com percentuais que variaram de 48,2 \% ( região Nordeste ) a 52,3 \% ( região Centro-Oeste ).
O percentual de alunos que citaram a Falta de motivação como dificuldade variou de 11,5 \% ( região Centro-Oeste ) a 16,1 \% ( região Nordeste ).
Os que declararam não ter qualquer tipo de dificuldade para responder à prova variaram de 18,1 \%, na região Nordeste, a 20,7 \%, na região Sul.
Em relação aos quartos de desempenho, o Desconhecimento do conteúdo foi a opção escolhida por 14,1 \% dos estudantes do quarto inferior e por 9,4 \% do quarto superior, caracterizando uma tendência decrescente.
São observadas diferenças estatisticamente significativas entre todos os quartos de desempenho.
A alternativa modal para os alunos, quando agregados pelos quartos de desempenho, para a dificuldade encontrada, foi causada pela Forma diferente de abordagem do conteúdo : 47,8 \%, no primeiro quarto, e 52,0 \% do terceiro quarto, escolheram essa alternativa.
Na análise por Organização Acadêmica, o percentual de inscritos e presentes que apontaram o Desconhecimento do conteúdo como “ dificuldade ao responder à prova ” não superou 13,0 \%.
Os percentuais variaram de 8,3 \%, nos CEFET/IF, a 13,0 \%, nos Centros Universitários.
A Forma diferente de abordagem do conteúdo foi a escolha modal dos estudantes, com percentuais que variaram de 49,1 \% ( Centros Universitários ) a 55,9 \% ( CEFET/IF ).
O percentual de alunos que citaram a Falta de motivação como dificuldade variou de 12,3 \% ( Faculdades ) a 14,8 \% ( CEFET/IF ).
Os que declararam não ter qualquer tipo de dificuldade para responder à prova variaram de 16,5 \%, nos CEFET/IF, a 20,2 \%, nas Universidades.
No Gráfico 4.27, é possível observar que há diferenças estatisticamente significativas entre os CEFET/IF e as demais Organizações Acadêmicas, bem como entre os Centros Universitários, Faculdades e Universidades ( ver também Tabela II.14, no Anexo II ).
Em relação às Categorias Administrativas, o Desconhecimento do conteúdo foi a opção escolhida por 8,6 \% dos estudantes, nas IES Públicas, e por 12,2 \%, nas IES Privadas.
São observadas diferenças estatisticamente significativas entre essas duas categorias.
A alternativa modal entre para os alunos, quando agregados pelas Categorias Administrativas, para a dificuldade encontrada, foi causada pela Forma diferente de abordagem do conteúdo : 48,0 \%, entre os estudantes de IES Públicas, e 50,3 \%, entre os de IES Privadas ( ver Gráfico 4.28 e Tabela II.14, no Anexo II ). 4.6 CONTEÚDOS DAS QUESTÕES OBJETIVAS DA PROVA Ao se analisarem os conteúdos das questões objetivas da prova ( Questão 8 ), verifica -se que um percentual muito pequeno dos estudantes avaliados, apenas 6,7 \%, afirmou que não estudou ainda a maioria desses conteúdos ( Gráficos 4.29 a Gráfico 4.32, e as Tabelas II.15 e II.16, no Anexo II ).
A maioria ( 62,6 \% ) afirmou ter estudado e aprendido muitos ou todos os conteúdos avaliados.
Na análise por Grande Região, a proporção de respondentes que escolheram a opção não estudou ainda a maioria desses conteúdos foi pequena.
Observa -se que, nas regiões Norte ( 6,9 \% ) e Sul ( 7,1 \% ), as proporções foram ligeiramente maiores do que a da média nacional ( 6,7 \% ).
Nas regiões Sudeste e Centro-Oeste, a situação foi inversa, com proporções mais baixas, correspondendo a, respectivamente, 6,4 \% e 6,1 \%.
Na região Nordeste, a proporção de respondentes que escolheram essa resposta foi a mesma da nacional ( 6,7 \% ).
São observadas diferenças estatisticamente significativas entre a região Sul e as regiões Sudeste e Centro-Oeste.
Em todas as regiões, a maioria dos presentes afirmou ter estudado e aprendido muitos ou todos os conteúdos, com proporções variando entre 61,1 \%, na região Centro-Oeste, a 63,9 \%, nas regiões Nordeste e Sudeste.
Considerando -se separadamente as opiniões de estudantes dos quatro quartos de desempenho, observa -se que, no quarto inferior, 12,4 \% ofereceram como resposta que não estudou ainda a maioria desses conteúdos, sendo 2,1 \% os do quarto superior com a mesma resposta.
A diferença entre os alunos que optaram por esse motivo de dificuldade entre todos os quartos de desempenho é estatisticamente significativa.
Tendo -se em conta o quarto superior, 80,7 \% dos alunos afirmaram ter estudado e aprendido muitos ou todos os conteúdos.
No outro extremo, no primeiro quarto, 45,7 \% optaram pelas mesmas categorias.
As proporções são crescentes com o desempenho.
Na análise por Organização Acadêmica, a proporção de respondentes que escolheram a opção não estudou ainda a maioria desses conteúdos foi pequena.
Observa -se que as proporções nos CEFET/IF ( 3,8 \% ) e nas Universidades ( 6,1 \% ) foram menores do que a da média nacional ( 6,7 \% ).
Nas Faculdades e nos Centros Universitários, a situação foi inversa : a proporção foi mais alta, 6,8 \% e 8,4 \%, respectivamente.
No Gráfico 4.31, é possível observar que as diferenças entre todos os tipos de Organização Acadêmica são estatisticamente significativas.
Em todos os tipos de Organização Acadêmica, a maioria dos presentes afirmou ter estudado e aprendido muitos ou todos os conteúdos, com proporções variando entre 60,2 \%, nos Centros Universitários, a 71,1 \%, nos CEFET/IF ( ver também Tabela II.16, no Anexo II ).
Em relação às Categorias Administrativas, observa -se que, nas IES Públicas, 3,4 \% ofereceram como resposta que não estudou ainda a maioria desses conteúdos, sendo 7,3 \% nas IES Privadas com a mesma resposta.
A diferença entre os alunos que optaram por esse motivo de dificuldade nas duas categorias é estatisticamente significativa.
Considerando -se as IES Públicas, 73,6 \% dos alunos afirmaram ter estudado e aprendido muitos ou todos os conteúdos, sendo de 60,5 \% a proporção dos que escolheram essa resposta nas IES Privadas ( ver Gráfico 4.32 e Tabela II.16, no Anexo II ). 4.7 TEMPO GASTO PARA CONCLUIR A PROVA Ao responderem sobre o tempo de conclusão da prova ( Questão 9 ), mais de quatro quintos dos estudantes ( 79,0 \% ) afirmaram ter gasto Entre duas e quatro horas ( Gráfico 4.33 a Gráfico 4.36 e, no Anexo II, as Tabelas II.17 e II.18 ).
Considerando -se as cinco Grandes Regiões brasileiras, os percentuais dos que utilizaram Entre duas e quatro horas para finalizar a prova nas regiões Norte ( 78,2 \% ) e Sudeste ( 78,6 \% ) foram inferiores ao percentual nacional.
Na região Nordeste, o percentual foi igual ao do total : 79,0 \%.
Nas regiões Sul e Centro-Oeste, os percentuais de alunos que dispensaram Entre duas e quatro horas para concluir a prova ficaram muito próximos e ligeiramente superiores à média : 79,1 \% e 80,4 \%, respectivamente.
Como mostrado no Gráfico 4.33, existe diferença estatisticamente significativa entre a região Centro-Oeste e as regiões Nordeste, Sudeste e Sul.
Uma vez analisadas as alternativas escolhidas pelos estudantes que se situam nos diferentes quartos de desempenho, observa -se uma tendência crescente : uma proporção maior de participantes no quarto superior declarou ter gasto Entre duas e quatro horas para concluir a prova quando comparada com a proporção de participantes nos quartos inferiores.
Há diferenças estatisticamente significativas entre todos os quartos de desempenho.
Na análise por Organização Acadêmica, os percentuais dos que utilizaram Entre duas e quatro horas para finalizar a prova nas Universidades ( 79,1 \% ), Faculdades ( 79,5 \% ) e CEFET/IF ( 83,4 \% ) foram superiores ao percentual nacional ( 79,0 \% ).
Nos Centros Universitários, tal percentual foi de 78,5 \%, e, como mostrado no Gráfico 4.35, há diferença estatisticamente significativa entre a proporção dos CEFET/IF e as proporções dos outros três tipos de Organização Acadêmica.
Em relação às Categorias Administrativas, observa -se que, nas IES Públicas, 79,6 \% declararam ter gasto Entre duas e quatro horas para concluir a prova, sendo 78,9 \% nas IES Privadas com a mesma resposta.
A diferença entre os alunos que optaram por esse motivo de dificuldade nas duas categorias não é estatisticamente significativa ( ver Gráfico 4.36 e Tabela II.18, no Anexo II ).

% matched lemmas: 
\end{textsample}
